\section{Resamplings on a vertex triple}\label{tri:section}

We now transfer the three-row inequality to graph resamplings whose
indexing pairs share a vertex. The internal edges of the three vertices
must be included in the matrix encoding; Table~\ref{tri:encoding}
records the four possible internal edge counts.
For $T\in\binom V3$, define
\[
 H_T=\sum_{a\in\binom T2}h_a.
\]
The three-row inequality yields the following local spectral bound.

\begin{proposition}\label{tri:gap}
For every vertex triple $T$, one has $H_T^2\succeq H_T$.
\end{proposition}
\begin{proof}
Partition $\Omega_d$ by recording the graph on $V\setminus T$ and,
for each $w\notin T$, its number of neighbors in $T$.
Every pair resampling indexed by a pair in $T$ preserves these data.
Work in one nonempty cell, with its conditional uniform measure.
The degree sum in $T$ determines the total number $e\in\{0,1,2,3\}$
of edges inside $T$.

Use a binary matrix with rows indexed by $i\in T$ and columns indexed
by $w\notin T$, whose entry at $(i,w)$ is the indicator of the edge
$iw$.  Its column sums are fixed.  To encode the internal graph as
well, let $s_i$ denote the degree of $i$ within $T$, and augment the
matrix as follows.
\begin{table}[ht]
\centering
\begin{tabular}{c|l|c|c}
$e$ & Internal graph & Extra column & Row sum at $i$\\ \hline
$0$ & Empty & None & $d_i$\\
$1$ & A single edge & $s_i$ (sum $2$) & $d_i$\\
$2$ & A path & $s_i-1$ (sum $1$) & $d_i-1$\\
$3$ & A triangle & None & $d_i-2$
\end{tabular}
\caption{Encoding the internal edges of a triple in a three-row matrix.
The extra column, when present, is a distinct labeled column.}
\label{tri:encoding}
\end{table}

For $e=1$, the extra column marks the endpoints of the sole edge;
for $e=2$, it marks the center of the path.  Thus the internal graph is
uniquely recovered from the extra column, or is deterministic when
$e=0,3$.  Conversely, every binary matrix with the prescribed row and
column sums recovers a valid graph in this cell.  The row sums give
the required degrees in $T$, while the external column sums and the
fixed graph on $V\setminus T$ give all other degrees.  This is a
bijection onto the entire fixed-margin matrix space, so its
conditional measure is uniform.

We next identify the conditional operators exactly.  Fix $i\in T$
and write $\{j,k\}=T\setminus\{i\}$.  Within this cell,
\begin{equation}\label{tri:conditional-row}
 E_{\{j,k\}}f=\mathbb E[f\mid S_i],
\end{equation}
where $S_i$ is augmented row $i$.
Indeed, the external entries of this row determine $s_i$ from $d_i$,
and hence determine its extra entry whenever present.  Fixing the
row therefore means precisely fixing the external neighbors of $i$.
Together with the cell data, these neighbors determine all the
$\{j,k\}$-fiber data:
\begin{itemize}
\item the graph away from $\{j,k\}$ is the fixed graph on
$V\setminus T$, together with the specified external edges at $i$;
\item the indicator of $jk$ is $e-s_i$;
\item the number of neighbors in $\{j,k\}$ of each $w\notin T$ is
its recorded number of neighbors in $T$ minus the indicator of $iw$;
\item the number of neighbors in $\{j,k\}$ of $i$ itself is $s_i$.
\end{itemize}
Conversely, the pair-fiber data fix all external edges at $i$, and
therefore fix $S_i$.  The fiber is contained in the cell because the
cell data are preserved by the pair resampling.  Thus the two
conditionings coincide on the full pair fiber, proving
\eqref{tri:conditional-row}.  In particular, when $s_i=1$, the choice
between $ij$ and $ik$ remains free whenever the other constraints
permit it; conditioning on the augmented row does not freeze that
internal choice.

Let $P_i f=\mathbb E[f\mid S_i]$.  For a function $f$ centered in the
cell, the functions $P_i f$ are centered row functions.  Put
$b=\sum_i\|P_i f\|^2$.  Orthogonality of $P_i$,
Lemma~\ref{row:variance}, and Cauchy--Schwarz yield
\[
 b=\left\langle f,\sum_i P_i f\right\rangle
 \le\|f\|\left\|\sum_i P_i f\right\|
 \le\|f\|\sqrt{2b}.
\]
It follows that $b\le2\|f\|^2$, including when $b=0$.
By \eqref{tri:conditional-row},
\[
 \langle f,H_Tf\rangle
 =3\|f\|^2-\sum_i\|P_i f\|^2\ge\|f\|^2.
\]
Thus the self-adjoint operator $H_T$ kills constants and has all its
other eigenvalues at least $1$ in this cell.  Squaring its eigenvalues
proves $H_T^2\succeq H_T$ there; the conclusion is also immediate for
a singleton cell.  Finally, all the operators preserve the cells, so
the cellwise inequalities give the claimed inequality on $\Omega_d$.
\end{proof}
